% VGG-T-Det with SAM 3 + VGGT-SLAM objects (task2-det): design discussion and the check made before any training.
% Fragment for % VGG-T-Det with SAM 3 + VGGT-SLAM objects (task2-det): design discussion and the check made before any training.
% Fragment for % VGG-T-Det with SAM 3 + VGGT-SLAM objects (task2-det): design discussion and the check made before any training.
% Fragment for % VGG-T-Det with SAM 3 + VGGT-SLAM objects (task2-det): design discussion and the check made before any training.
% Fragment for \input{vggt-det-arch}; needs \usepackage{booktabs}. Refers to sam3.tex (\ref{sec:sam3}, tab:sam3-names).
% Status: design and measurements only; no detector has been trained with these changes yet.
% Check: VGGT-SLAM/exps/vsibench_objects/query_cov2.py on the 88 ScanNet VSI-Bench scenes,
%        ground truth vggtdet_es56/data/scannet/scannet_infos_vsibench88_esl31.pkl (1,273 boxes, 16 classes),
%        SAM 3 lists all/s3pbox_<scene>.json (fixed 65-name prompts, plain ``door'').
% Detector code read: vggtdet_es56/code/projects/VGGTDet/vggtdet/{vggtdet.py, vggt_head.py, multiview_pipeline.py}.

\section{Using SAM 3 + VGGT-SLAM objects inside VGG-T-Det}
\label{sec:vggtdet-arch}

\paragraph{Motivation.}
The object list from SAM 3 and VGGT-SLAM (Section~\ref{sec:sam3}) has usable labels but wrong counts: extra
boxes and missing objects, not box geometry, limit its benefit, and box size is the only geometric quantity
whose correction helps. VGG-T-Det is trained to regress boxes and to reject false ones, but its queries miss
small objects. We therefore ask how the detector can start from
the SAM 3 objects: they propose where objects are and what they are called; the detector refines the box and
decides which proposals to keep.

\paragraph{The detector today.}
VGG-T-Det takes 256 queries by farthest-point sampling from the point cloud VGGT predicts for the input frames.
A query carries only the position embedding of its point; its content vector starts at zero. A transformer
decoder attends from the queries to the VGGT tokens, and a head predicts for each query a centre offset, a size
and class scores. Predictions are assigned to ground-truth boxes by Hungarian matching on class score and
generalised IoU, with additional matches for the classes the original model did not have.
An earlier memorisation test on eight scenes showed that localisation, not classification, fails for small
objects: a prediction that overlaps a ground-truth box has the right label in 97\% of cases, but only 27\% of the
boxes had a query inside them and VGGT's cloud had no point at all inside 32\% of them.

\subsection{Queries from the detected objects (task2-det-query)}
\label{sec:vggtdet-arch-query}

\paragraph{Design.}
Each SAM 3 + VGGT-SLAM box contributes queries drawn from its own object points, so that an object has queries
however small it is:
\begin{itemize}
\item at least five queries per box, and more for large boxes:
      $n = \min(12,\; 5 + \mathrm{round}(\text{largest side} / 0.5\,\mathrm{m}))$;
\item the first query is the object point nearest to the box centre; the others are drawn with probability
      proportional to the local point density times a weight that decreases with the distance from the centre;
\item the 256 farthest-point queries are kept, and the object queries are added to them.
\end{itemize}
The detector already has a switch that places given 3D points as queries (built for a test with ground-truth
centres), so the object queries can enter through the same path.

\paragraph{Frame and scale.}
The SLAM points are in the SLAM frame and scale; the queries live in the frame of the first input image.
The detector, as it is run today, receives the ground-truth pose of the first frame and a ground-truth scale
(the mean distance of the scan points from the first camera) in training and at test time. Aligning the SLAM
objects to the scene with the true scale, by a similarity fit of the camera centres, is therefore consistent
with the present protocol. A detector that sees the video only would have to use the list's own scale estimate
instead; Table~\ref{tab:arch-cov} shows what that costs.

\paragraph{Check before training: do the queries land in the ground-truth boxes?}
We sampled queries as above for the 88 ScanNet scenes of VSI-Bench and counted the ground-truth boxes of the
detector that contain at least one query (Table~\ref{tab:arch-cov}).

\begin{table}[t]
\centering
\caption{Ground-truth boxes that contain at least one query (88 scenes, 1,273 boxes), by the box's largest side.
Farthest-point: 256 samples of the SLAM cloud. Object queries: from SAM 3 boxes above the given score.
Centre: boxes with a query within 0.1\,m of the box centre.}
\label{tab:arch-cov}
\begin{tabular}{lrrrrrrr}
\toprule
Query set & Queries & All & $<0.3$\,m & 0.3--0.6\,m & 0.6--1.2\,m & $>1.2$\,m & Centre \\
\midrule
\multicolumn{8}{l}{\emph{Aligned with the true scale}} \\
Farthest-point, 256              & 256 & 84.0\% &  3.7\% & 65.8\% & 85.1\% & 93.8\% &  4.8\% \\
Object queries, no threshold     & 362 & 84.3\% & 25.9\% & 80.0\% & 85.5\% & 87.7\% & 16.9\% \\
Object queries, score $\geq 0.3$ & 302 & 82.4\% & 22.2\% & 78.9\% & 83.9\% & 85.3\% & 14.5\% \\
Object queries, score $\geq 0.7$ & 188 & 77.5\% & 14.8\% & 74.2\% & 79.2\% & 80.2\% & 13.6\% \\
Object ($\geq 0.3$) + farthest-point, 256 in total & 256 & 78.2\% & 14.8\% & 73.7\% & 79.2\% & 82.3\% & 13.2\% \\
\midrule
\multicolumn{8}{l}{\emph{Aligned with the list's own scale estimate}} \\
Farthest-point, 256              & 256 & 70.3\% &  0.0\% & 52.1\% & 74.5\% & 76.4\% &  2.4\% \\
Object queries, score $\geq 0.3$ & 302 & 67.6\% & 14.8\% & 61.6\% & 73.0\% & 66.9\% &  9.9\% \\
\bottomrule
\end{tabular}
\end{table}

\begin{itemize}
\item \emph{Small objects gain.} Below 0.6\,m the object queries cover 79--80\% of the boxes against 66\% for
      farthest-point sampling; below 0.3\,m, 22--26\% against 4\% (27 boxes only).
\item \emph{Large objects lose.} Above 1.2\,m farthest-point sampling covers 94\%, object queries 85--88\%,
      because SAM 3 misses some large objects entirely.
\item \emph{Centres are closer.} Three times as many boxes have a query within 0.1\,m of their centre.
\item \emph{Replacing is worse than adding.} With a fixed budget of 256, mixing the two kinds covers fewer boxes
      (78\%) than either kind alone. The object queries must be added to the 256, not take their place.
\item \emph{Scale matters more than the kind of query.} With the estimated scale every row loses about 14
      points.
\end{itemize}

\subsection{Labels from the proposals and a keep decision (task2-det-matching)}
\label{sec:vggtdet-arch-match}

\paragraph{Basic test: the detector as a refiner.}
The label of an object query is the label of its SAM 3 box, and the class head is not used. A new output decides
for each SAM 3 box whether it is kept.
\begin{itemize}
\item Each SAM 3 box is assigned beforehand to at most one ground-truth box of the same label (its centre inside
      the ground-truth box, or the reverse, or IoU $\geq 0.25$; one-to-one, best score first).
\item The queries of an assigned box regress the centre and size of that ground-truth box and have keep $=1$.
      The queries of an unassigned box have keep $=0$ and no box loss. No Hungarian matching is needed for them.
\item At test time the queries of one SAM 3 box are pooled into one output box, so no non-maximum suppression is
      needed.
\end{itemize}
The detector then only refines centre and size and removes extra boxes, which is what the list lacks.

\paragraph{Ceiling of the basic test.}
Table~\ref{tab:arch-keep} gives the result of the assignment on the 88 scenes, for the classes that have
ground truth there.

\begin{table}[t]
\centering
\caption{SAM 3 boxes assigned one-to-one to ground-truth boxes of the same label (88 scenes, true scale).}
\label{tab:arch-keep}
\begin{tabular}{rrrrr}
\toprule
SAM 3 score & Boxes per scene & Keep & Drop & Ground truth with a same-label box \\
\midrule
none       & 31.9 & 31.0\% & 69.0\% & 68.4\% \\
$\geq 0.3$ & 26.0 & 37.2\% & 62.8\% & 66.8\% \\
$\geq 0.7$ & 17.0 & 52.9\% & 47.1\% & 62.4\% \\
\bottomrule
\end{tabular}
\end{table}

\begin{itemize}
\item \emph{Recall is capped at about two thirds.} A detector that only refines and filters cannot find the
      remaining third, and a lower threshold changes little (62\% to 68\%).
\item \emph{The weak classes} are picture (20\% of its boxes contain a same-label query at score $\geq 0.3$),
      trash bin (48\%), cabinet (57\%), window (59\%) and door (67\%); chair, sofa, bed, toilet and desk are at
      89--96\%.
\item \emph{The keep target is not rare.} Between half and two thirds of the proposals are negatives, so the keep
      output has examples of both kinds.
\end{itemize}

\paragraph{Two limits of the basic test.}
It cannot add an object that SAM 3 missed. And a correct SAM 3 box on an object that has no annotation is trained
as ``drop''; this is most likely for ARKitScenes, so the test starts on ScanNet.

\paragraph{Second stage: both kinds of queries.}
The free queries keep the class head and Hungarian matching and can find the objects SAM 3 missed. For the object
queries the label enters in two places: in the matching cost, as a penalty for a ground-truth box of another
label or as a restriction to boxes of the same label; and in the query content, as an embedding of the label
together with the SAM 3 score and the proposal's size, where the content is zero today.

\subsection{Object features from the 2D detections (task2-det-feature)}
\label{sec:vggtdet-arch-feat}

A SAM 3 object is seen as a mask in several frames. Its 2D evidence can be aggregated over those frames into
one feature vector and given to the detector as the content of the object's queries. Three sources are possible
(Table~\ref{tab:arch-feat}).

\begin{table}[t]
\centering
\caption{Sources for a per-object feature, cheapest first.}
\label{tab:arch-feat}
\begin{tabular}{llp{0.42\linewidth}}
\toprule
 & Source & Cost \\
\midrule
A & The detector's own VGGT tokens, averaged under each mask and over the object's views, weighted by the
    SAM 3 score & No new SAM 3 run; the detector must be given the SLAM keyframes as input frames \\
B & SAM 3's own embedding of each detection, averaged over the views & SAM 3 must be run again to store them
    (only masks and scores are stored now) \\
C & SAM 3 image-encoder features pooled under the mask & Run again, and larger files \\
\bottomrule
\end{tabular}
\end{table}

Option A comes first: it needs no new data and tells the decoder which image regions belong to the object.
Option B carries the strongest semantic signal and is added only if A helps.

\subsection{Order of work}

\begin{enumerate}
\item Query coverage and the ceiling of the keep decision, without training
      (done on 88 scenes: Tables~\ref{tab:arch-cov} and~\ref{tab:arch-keep}).
\item The same measurement on the ScanNet training scenes, whose 65-class ground truth includes small objects.
\item Basic test on ScanNet: object queries added to the 256 free queries, label from the proposal, keep output,
      regression targets from the fixed assignment. Validation on the 88 ScanNet scenes of VSI-Bench.
\item Second stage with both kinds of queries.
\item Object features, option A.
\end{enumerate}
Each stage is judged by the list counts (extra and missing boxes) and by the VSI-Bench score of the language
model against the unrefined SAM 3 list.

\paragraph{Limits of the check so far.}
The farthest-point row samples the SLAM cloud as a stand-in for the detector's queries, which are sampled from
VGGT's own cloud for the input frames. The ground truth of the 88 scenes has boxes for 16 classes of large
furniture only, so small objects are under-represented. The SAM 3 lists used the fixed 65-name prompts with
plain ``door'', not ``room door''.
; needs \usepackage{booktabs}. Refers to sam3.tex (\ref{sec:sam3}, tab:sam3-names).
% Status: design and measurements only; no detector has been trained with these changes yet.
% Check: VGGT-SLAM/exps/vsibench_objects/query_cov2.py on the 88 ScanNet VSI-Bench scenes,
%        ground truth vggtdet_es56/data/scannet/scannet_infos_vsibench88_esl31.pkl (1,273 boxes, 16 classes),
%        SAM 3 lists all/s3pbox_<scene>.json (fixed 65-name prompts, plain ``door'').
% Detector code read: vggtdet_es56/code/projects/VGGTDet/vggtdet/{vggtdet.py, vggt_head.py, multiview_pipeline.py}.

\section{Using SAM 3 + VGGT-SLAM objects inside VGG-T-Det}
\label{sec:vggtdet-arch}

\paragraph{Motivation.}
The object list from SAM 3 and VGGT-SLAM (Section~\ref{sec:sam3}) has usable labels but wrong counts: extra
boxes and missing objects, not box geometry, limit its benefit, and box size is the only geometric quantity
whose correction helps. VGG-T-Det is trained to regress boxes and to reject false ones, but its queries miss
small objects. We therefore ask how the detector can start from
the SAM 3 objects: they propose where objects are and what they are called; the detector refines the box and
decides which proposals to keep.

\paragraph{The detector today.}
VGG-T-Det takes 256 queries by farthest-point sampling from the point cloud VGGT predicts for the input frames.
A query carries only the position embedding of its point; its content vector starts at zero. A transformer
decoder attends from the queries to the VGGT tokens, and a head predicts for each query a centre offset, a size
and class scores. Predictions are assigned to ground-truth boxes by Hungarian matching on class score and
generalised IoU, with additional matches for the classes the original model did not have.
An earlier memorisation test on eight scenes showed that localisation, not classification, fails for small
objects: a prediction that overlaps a ground-truth box has the right label in 97\% of cases, but only 27\% of the
boxes had a query inside them and VGGT's cloud had no point at all inside 32\% of them.

\subsection{Queries from the detected objects (task2-det-query)}
\label{sec:vggtdet-arch-query}

\paragraph{Design.}
Each SAM 3 + VGGT-SLAM box contributes queries drawn from its own object points, so that an object has queries
however small it is:
\begin{itemize}
\item at least five queries per box, and more for large boxes:
      $n = \min(12,\; 5 + \mathrm{round}(\text{largest side} / 0.5\,\mathrm{m}))$;
\item the first query is the object point nearest to the box centre; the others are drawn with probability
      proportional to the local point density times a weight that decreases with the distance from the centre;
\item the 256 farthest-point queries are kept, and the object queries are added to them.
\end{itemize}
The detector already has a switch that places given 3D points as queries (built for a test with ground-truth
centres), so the object queries can enter through the same path.

\paragraph{Frame and scale.}
The SLAM points are in the SLAM frame and scale; the queries live in the frame of the first input image.
The detector, as it is run today, receives the ground-truth pose of the first frame and a ground-truth scale
(the mean distance of the scan points from the first camera) in training and at test time. Aligning the SLAM
objects to the scene with the true scale, by a similarity fit of the camera centres, is therefore consistent
with the present protocol. A detector that sees the video only would have to use the list's own scale estimate
instead; Table~\ref{tab:arch-cov} shows what that costs.

\paragraph{Check before training: do the queries land in the ground-truth boxes?}
We sampled queries as above for the 88 ScanNet scenes of VSI-Bench and counted the ground-truth boxes of the
detector that contain at least one query (Table~\ref{tab:arch-cov}).

\begin{table}[t]
\centering
\caption{Ground-truth boxes that contain at least one query (88 scenes, 1,273 boxes), by the box's largest side.
Farthest-point: 256 samples of the SLAM cloud. Object queries: from SAM 3 boxes above the given score.
Centre: boxes with a query within 0.1\,m of the box centre.}
\label{tab:arch-cov}
\begin{tabular}{lrrrrrrr}
\toprule
Query set & Queries & All & $<0.3$\,m & 0.3--0.6\,m & 0.6--1.2\,m & $>1.2$\,m & Centre \\
\midrule
\multicolumn{8}{l}{\emph{Aligned with the true scale}} \\
Farthest-point, 256              & 256 & 84.0\% &  3.7\% & 65.8\% & 85.1\% & 93.8\% &  4.8\% \\
Object queries, no threshold     & 362 & 84.3\% & 25.9\% & 80.0\% & 85.5\% & 87.7\% & 16.9\% \\
Object queries, score $\geq 0.3$ & 302 & 82.4\% & 22.2\% & 78.9\% & 83.9\% & 85.3\% & 14.5\% \\
Object queries, score $\geq 0.7$ & 188 & 77.5\% & 14.8\% & 74.2\% & 79.2\% & 80.2\% & 13.6\% \\
Object ($\geq 0.3$) + farthest-point, 256 in total & 256 & 78.2\% & 14.8\% & 73.7\% & 79.2\% & 82.3\% & 13.2\% \\
\midrule
\multicolumn{8}{l}{\emph{Aligned with the list's own scale estimate}} \\
Farthest-point, 256              & 256 & 70.3\% &  0.0\% & 52.1\% & 74.5\% & 76.4\% &  2.4\% \\
Object queries, score $\geq 0.3$ & 302 & 67.6\% & 14.8\% & 61.6\% & 73.0\% & 66.9\% &  9.9\% \\
\bottomrule
\end{tabular}
\end{table}

\begin{itemize}
\item \emph{Small objects gain.} Below 0.6\,m the object queries cover 79--80\% of the boxes against 66\% for
      farthest-point sampling; below 0.3\,m, 22--26\% against 4\% (27 boxes only).
\item \emph{Large objects lose.} Above 1.2\,m farthest-point sampling covers 94\%, object queries 85--88\%,
      because SAM 3 misses some large objects entirely.
\item \emph{Centres are closer.} Three times as many boxes have a query within 0.1\,m of their centre.
\item \emph{Replacing is worse than adding.} With a fixed budget of 256, mixing the two kinds covers fewer boxes
      (78\%) than either kind alone. The object queries must be added to the 256, not take their place.
\item \emph{Scale matters more than the kind of query.} With the estimated scale every row loses about 14
      points.
\end{itemize}

\subsection{Labels from the proposals and a keep decision (task2-det-matching)}
\label{sec:vggtdet-arch-match}

\paragraph{Basic test: the detector as a refiner.}
The label of an object query is the label of its SAM 3 box, and the class head is not used. A new output decides
for each SAM 3 box whether it is kept.
\begin{itemize}
\item Each SAM 3 box is assigned beforehand to at most one ground-truth box of the same label (its centre inside
      the ground-truth box, or the reverse, or IoU $\geq 0.25$; one-to-one, best score first).
\item The queries of an assigned box regress the centre and size of that ground-truth box and have keep $=1$.
      The queries of an unassigned box have keep $=0$ and no box loss. No Hungarian matching is needed for them.
\item At test time the queries of one SAM 3 box are pooled into one output box, so no non-maximum suppression is
      needed.
\end{itemize}
The detector then only refines centre and size and removes extra boxes, which is what the list lacks.

\paragraph{Ceiling of the basic test.}
Table~\ref{tab:arch-keep} gives the result of the assignment on the 88 scenes, for the classes that have
ground truth there.

\begin{table}[t]
\centering
\caption{SAM 3 boxes assigned one-to-one to ground-truth boxes of the same label (88 scenes, true scale).}
\label{tab:arch-keep}
\begin{tabular}{rrrrr}
\toprule
SAM 3 score & Boxes per scene & Keep & Drop & Ground truth with a same-label box \\
\midrule
none       & 31.9 & 31.0\% & 69.0\% & 68.4\% \\
$\geq 0.3$ & 26.0 & 37.2\% & 62.8\% & 66.8\% \\
$\geq 0.7$ & 17.0 & 52.9\% & 47.1\% & 62.4\% \\
\bottomrule
\end{tabular}
\end{table}

\begin{itemize}
\item \emph{Recall is capped at about two thirds.} A detector that only refines and filters cannot find the
      remaining third, and a lower threshold changes little (62\% to 68\%).
\item \emph{The weak classes} are picture (20\% of its boxes contain a same-label query at score $\geq 0.3$),
      trash bin (48\%), cabinet (57\%), window (59\%) and door (67\%); chair, sofa, bed, toilet and desk are at
      89--96\%.
\item \emph{The keep target is not rare.} Between half and two thirds of the proposals are negatives, so the keep
      output has examples of both kinds.
\end{itemize}

\paragraph{Two limits of the basic test.}
It cannot add an object that SAM 3 missed. And a correct SAM 3 box on an object that has no annotation is trained
as ``drop''; this is most likely for ARKitScenes, so the test starts on ScanNet.

\paragraph{Second stage: both kinds of queries.}
The free queries keep the class head and Hungarian matching and can find the objects SAM 3 missed. For the object
queries the label enters in two places: in the matching cost, as a penalty for a ground-truth box of another
label or as a restriction to boxes of the same label; and in the query content, as an embedding of the label
together with the SAM 3 score and the proposal's size, where the content is zero today.

\subsection{Object features from the 2D detections (task2-det-feature)}
\label{sec:vggtdet-arch-feat}

A SAM 3 object is seen as a mask in several frames. Its 2D evidence can be aggregated over those frames into
one feature vector and given to the detector as the content of the object's queries. Three sources are possible
(Table~\ref{tab:arch-feat}).

\begin{table}[t]
\centering
\caption{Sources for a per-object feature, cheapest first.}
\label{tab:arch-feat}
\begin{tabular}{llp{0.42\linewidth}}
\toprule
 & Source & Cost \\
\midrule
A & The detector's own VGGT tokens, averaged under each mask and over the object's views, weighted by the
    SAM 3 score & No new SAM 3 run; the detector must be given the SLAM keyframes as input frames \\
B & SAM 3's own embedding of each detection, averaged over the views & SAM 3 must be run again to store them
    (only masks and scores are stored now) \\
C & SAM 3 image-encoder features pooled under the mask & Run again, and larger files \\
\bottomrule
\end{tabular}
\end{table}

Option A comes first: it needs no new data and tells the decoder which image regions belong to the object.
Option B carries the strongest semantic signal and is added only if A helps.

\subsection{Order of work}

\begin{enumerate}
\item Query coverage and the ceiling of the keep decision, without training
      (done on 88 scenes: Tables~\ref{tab:arch-cov} and~\ref{tab:arch-keep}).
\item The same measurement on the ScanNet training scenes, whose 65-class ground truth includes small objects.
\item Basic test on ScanNet: object queries added to the 256 free queries, label from the proposal, keep output,
      regression targets from the fixed assignment. Validation on the 88 ScanNet scenes of VSI-Bench.
\item Second stage with both kinds of queries.
\item Object features, option A.
\end{enumerate}
Each stage is judged by the list counts (extra and missing boxes) and by the VSI-Bench score of the language
model against the unrefined SAM 3 list.

\paragraph{Limits of the check so far.}
The farthest-point row samples the SLAM cloud as a stand-in for the detector's queries, which are sampled from
VGGT's own cloud for the input frames. The ground truth of the 88 scenes has boxes for 16 classes of large
furniture only, so small objects are under-represented. The SAM 3 lists used the fixed 65-name prompts with
plain ``door'', not ``room door''.
; needs \usepackage{booktabs}. Refers to sam3.tex (\ref{sec:sam3}, tab:sam3-names).
% Status: design and measurements only; no detector has been trained with these changes yet.
% Check: VGGT-SLAM/exps/vsibench_objects/query_cov2.py on the 88 ScanNet VSI-Bench scenes,
%        ground truth vggtdet_es56/data/scannet/scannet_infos_vsibench88_esl31.pkl (1,273 boxes, 16 classes),
%        SAM 3 lists all/s3pbox_<scene>.json (fixed 65-name prompts, plain ``door'').
% Detector code read: vggtdet_es56/code/projects/VGGTDet/vggtdet/{vggtdet.py, vggt_head.py, multiview_pipeline.py}.

\section{Using SAM 3 + VGGT-SLAM objects inside VGG-T-Det}
\label{sec:vggtdet-arch}

\paragraph{Motivation.}
The object list from SAM 3 and VGGT-SLAM (Section~\ref{sec:sam3}) has usable labels but wrong counts: extra
boxes and missing objects, not box geometry, limit its benefit, and box size is the only geometric quantity
whose correction helps. VGG-T-Det is trained to regress boxes and to reject false ones, but its queries miss
small objects. We therefore ask how the detector can start from
the SAM 3 objects: they propose where objects are and what they are called; the detector refines the box and
decides which proposals to keep.

\paragraph{The detector today.}
VGG-T-Det takes 256 queries by farthest-point sampling from the point cloud VGGT predicts for the input frames.
A query carries only the position embedding of its point; its content vector starts at zero. A transformer
decoder attends from the queries to the VGGT tokens, and a head predicts for each query a centre offset, a size
and class scores. Predictions are assigned to ground-truth boxes by Hungarian matching on class score and
generalised IoU, with additional matches for the classes the original model did not have.
An earlier memorisation test on eight scenes showed that localisation, not classification, fails for small
objects: a prediction that overlaps a ground-truth box has the right label in 97\% of cases, but only 27\% of the
boxes had a query inside them and VGGT's cloud had no point at all inside 32\% of them.

\subsection{Queries from the detected objects (task2-det-query)}
\label{sec:vggtdet-arch-query}

\paragraph{Design.}
Each SAM 3 + VGGT-SLAM box contributes queries drawn from its own object points, so that an object has queries
however small it is:
\begin{itemize}
\item at least five queries per box, and more for large boxes:
      $n = \min(12,\; 5 + \mathrm{round}(\text{largest side} / 0.5\,\mathrm{m}))$;
\item the first query is the object point nearest to the box centre; the others are drawn with probability
      proportional to the local point density times a weight that decreases with the distance from the centre;
\item the 256 farthest-point queries are kept, and the object queries are added to them.
\end{itemize}
The detector already has a switch that places given 3D points as queries (built for a test with ground-truth
centres), so the object queries can enter through the same path.

\paragraph{Frame and scale.}
The SLAM points are in the SLAM frame and scale; the queries live in the frame of the first input image.
The detector, as it is run today, receives the ground-truth pose of the first frame and a ground-truth scale
(the mean distance of the scan points from the first camera) in training and at test time. Aligning the SLAM
objects to the scene with the true scale, by a similarity fit of the camera centres, is therefore consistent
with the present protocol. A detector that sees the video only would have to use the list's own scale estimate
instead; Table~\ref{tab:arch-cov} shows what that costs.

\paragraph{Check before training: do the queries land in the ground-truth boxes?}
We sampled queries as above for the 88 ScanNet scenes of VSI-Bench and counted the ground-truth boxes of the
detector that contain at least one query (Table~\ref{tab:arch-cov}).

\begin{table}[t]
\centering
\caption{Ground-truth boxes that contain at least one query (88 scenes, 1,273 boxes), by the box's largest side.
Farthest-point: 256 samples of the SLAM cloud. Object queries: from SAM 3 boxes above the given score.
Centre: boxes with a query within 0.1\,m of the box centre.}
\label{tab:arch-cov}
\begin{tabular}{lrrrrrrr}
\toprule
Query set & Queries & All & $<0.3$\,m & 0.3--0.6\,m & 0.6--1.2\,m & $>1.2$\,m & Centre \\
\midrule
\multicolumn{8}{l}{\emph{Aligned with the true scale}} \\
Farthest-point, 256              & 256 & 84.0\% &  3.7\% & 65.8\% & 85.1\% & 93.8\% &  4.8\% \\
Object queries, no threshold     & 362 & 84.3\% & 25.9\% & 80.0\% & 85.5\% & 87.7\% & 16.9\% \\
Object queries, score $\geq 0.3$ & 302 & 82.4\% & 22.2\% & 78.9\% & 83.9\% & 85.3\% & 14.5\% \\
Object queries, score $\geq 0.7$ & 188 & 77.5\% & 14.8\% & 74.2\% & 79.2\% & 80.2\% & 13.6\% \\
Object ($\geq 0.3$) + farthest-point, 256 in total & 256 & 78.2\% & 14.8\% & 73.7\% & 79.2\% & 82.3\% & 13.2\% \\
\midrule
\multicolumn{8}{l}{\emph{Aligned with the list's own scale estimate}} \\
Farthest-point, 256              & 256 & 70.3\% &  0.0\% & 52.1\% & 74.5\% & 76.4\% &  2.4\% \\
Object queries, score $\geq 0.3$ & 302 & 67.6\% & 14.8\% & 61.6\% & 73.0\% & 66.9\% &  9.9\% \\
\bottomrule
\end{tabular}
\end{table}

\begin{itemize}
\item \emph{Small objects gain.} Below 0.6\,m the object queries cover 79--80\% of the boxes against 66\% for
      farthest-point sampling; below 0.3\,m, 22--26\% against 4\% (27 boxes only).
\item \emph{Large objects lose.} Above 1.2\,m farthest-point sampling covers 94\%, object queries 85--88\%,
      because SAM 3 misses some large objects entirely.
\item \emph{Centres are closer.} Three times as many boxes have a query within 0.1\,m of their centre.
\item \emph{Replacing is worse than adding.} With a fixed budget of 256, mixing the two kinds covers fewer boxes
      (78\%) than either kind alone. The object queries must be added to the 256, not take their place.
\item \emph{Scale matters more than the kind of query.} With the estimated scale every row loses about 14
      points.
\end{itemize}

\subsection{Labels from the proposals and a keep decision (task2-det-matching)}
\label{sec:vggtdet-arch-match}

\paragraph{Basic test: the detector as a refiner.}
The label of an object query is the label of its SAM 3 box, and the class head is not used. A new output decides
for each SAM 3 box whether it is kept.
\begin{itemize}
\item Each SAM 3 box is assigned beforehand to at most one ground-truth box of the same label (its centre inside
      the ground-truth box, or the reverse, or IoU $\geq 0.25$; one-to-one, best score first).
\item The queries of an assigned box regress the centre and size of that ground-truth box and have keep $=1$.
      The queries of an unassigned box have keep $=0$ and no box loss. No Hungarian matching is needed for them.
\item At test time the queries of one SAM 3 box are pooled into one output box, so no non-maximum suppression is
      needed.
\end{itemize}
The detector then only refines centre and size and removes extra boxes, which is what the list lacks.

\paragraph{Ceiling of the basic test.}
Table~\ref{tab:arch-keep} gives the result of the assignment on the 88 scenes, for the classes that have
ground truth there.

\begin{table}[t]
\centering
\caption{SAM 3 boxes assigned one-to-one to ground-truth boxes of the same label (88 scenes, true scale).}
\label{tab:arch-keep}
\begin{tabular}{rrrrr}
\toprule
SAM 3 score & Boxes per scene & Keep & Drop & Ground truth with a same-label box \\
\midrule
none       & 31.9 & 31.0\% & 69.0\% & 68.4\% \\
$\geq 0.3$ & 26.0 & 37.2\% & 62.8\% & 66.8\% \\
$\geq 0.7$ & 17.0 & 52.9\% & 47.1\% & 62.4\% \\
\bottomrule
\end{tabular}
\end{table}

\begin{itemize}
\item \emph{Recall is capped at about two thirds.} A detector that only refines and filters cannot find the
      remaining third, and a lower threshold changes little (62\% to 68\%).
\item \emph{The weak classes} are picture (20\% of its boxes contain a same-label query at score $\geq 0.3$),
      trash bin (48\%), cabinet (57\%), window (59\%) and door (67\%); chair, sofa, bed, toilet and desk are at
      89--96\%.
\item \emph{The keep target is not rare.} Between half and two thirds of the proposals are negatives, so the keep
      output has examples of both kinds.
\end{itemize}

\paragraph{Two limits of the basic test.}
It cannot add an object that SAM 3 missed. And a correct SAM 3 box on an object that has no annotation is trained
as ``drop''; this is most likely for ARKitScenes, so the test starts on ScanNet.

\paragraph{Second stage: both kinds of queries.}
The free queries keep the class head and Hungarian matching and can find the objects SAM 3 missed. For the object
queries the label enters in two places: in the matching cost, as a penalty for a ground-truth box of another
label or as a restriction to boxes of the same label; and in the query content, as an embedding of the label
together with the SAM 3 score and the proposal's size, where the content is zero today.

\subsection{Object features from the 2D detections (task2-det-feature)}
\label{sec:vggtdet-arch-feat}

A SAM 3 object is seen as a mask in several frames. Its 2D evidence can be aggregated over those frames into
one feature vector and given to the detector as the content of the object's queries. Three sources are possible
(Table~\ref{tab:arch-feat}).

\begin{table}[t]
\centering
\caption{Sources for a per-object feature, cheapest first.}
\label{tab:arch-feat}
\begin{tabular}{llp{0.42\linewidth}}
\toprule
 & Source & Cost \\
\midrule
A & The detector's own VGGT tokens, averaged under each mask and over the object's views, weighted by the
    SAM 3 score & No new SAM 3 run; the detector must be given the SLAM keyframes as input frames \\
B & SAM 3's own embedding of each detection, averaged over the views & SAM 3 must be run again to store them
    (only masks and scores are stored now) \\
C & SAM 3 image-encoder features pooled under the mask & Run again, and larger files \\
\bottomrule
\end{tabular}
\end{table}

Option A comes first: it needs no new data and tells the decoder which image regions belong to the object.
Option B carries the strongest semantic signal and is added only if A helps.

\subsection{Order of work}

\begin{enumerate}
\item Query coverage and the ceiling of the keep decision, without training
      (done on 88 scenes: Tables~\ref{tab:arch-cov} and~\ref{tab:arch-keep}).
\item The same measurement on the ScanNet training scenes, whose 65-class ground truth includes small objects.
\item Basic test on ScanNet: object queries added to the 256 free queries, label from the proposal, keep output,
      regression targets from the fixed assignment. Validation on the 88 ScanNet scenes of VSI-Bench.
\item Second stage with both kinds of queries.
\item Object features, option A.
\end{enumerate}
Each stage is judged by the list counts (extra and missing boxes) and by the VSI-Bench score of the language
model against the unrefined SAM 3 list.

\paragraph{Limits of the check so far.}
The farthest-point row samples the SLAM cloud as a stand-in for the detector's queries, which are sampled from
VGGT's own cloud for the input frames. The ground truth of the 88 scenes has boxes for 16 classes of large
furniture only, so small objects are under-represented. The SAM 3 lists used the fixed 65-name prompts with
plain ``door'', not ``room door''.
; needs \usepackage{booktabs}. Refers to sam3.tex (\ref{sec:sam3}, tab:sam3-names).
% Status: design and measurements only; no detector has been trained with these changes yet.
% Check: VGGT-SLAM/exps/vsibench_objects/query_cov2.py on the 88 ScanNet VSI-Bench scenes,
%        ground truth vggtdet_es56/data/scannet/scannet_infos_vsibench88_esl31.pkl (1,273 boxes, 16 classes),
%        SAM 3 lists all/s3pbox_<scene>.json (fixed 65-name prompts, plain ``door'').
% Detector code read: vggtdet_es56/code/projects/VGGTDet/vggtdet/{vggtdet.py, vggt_head.py, multiview_pipeline.py}.

\section{Using SAM 3 + VGGT-SLAM objects inside VGG-T-Det}
\label{sec:vggtdet-arch}

\paragraph{Motivation.}
The object list from SAM 3 and VGGT-SLAM (Section~\ref{sec:sam3}) has usable labels but wrong counts: extra
boxes and missing objects, not box geometry, limit its benefit, and box size is the only geometric quantity
whose correction helps. VGG-T-Det is trained to regress boxes and to reject false ones, but its queries miss
small objects. We therefore ask how the detector can start from
the SAM 3 objects: they propose where objects are and what they are called; the detector refines the box and
decides which proposals to keep.

\paragraph{The detector today.}
VGG-T-Det takes 256 queries by farthest-point sampling from the point cloud VGGT predicts for the input frames.
A query carries only the position embedding of its point; its content vector starts at zero. A transformer
decoder attends from the queries to the VGGT tokens, and a head predicts for each query a centre offset, a size
and class scores. Predictions are assigned to ground-truth boxes by Hungarian matching on class score and
generalised IoU, with additional matches for the classes the original model did not have.
An earlier memorisation test on eight scenes showed that localisation, not classification, fails for small
objects: a prediction that overlaps a ground-truth box has the right label in 97\% of cases, but only 27\% of the
boxes had a query inside them and VGGT's cloud had no point at all inside 32\% of them.

\subsection{Queries from the detected objects (task2-det-query)}
\label{sec:vggtdet-arch-query}

\paragraph{Design.}
Each SAM 3 + VGGT-SLAM box contributes queries drawn from its own object points, so that an object has queries
however small it is:
\begin{itemize}
\item at least five queries per box, and more for large boxes:
      $n = \min(12,\; 5 + \mathrm{round}(\text{largest side} / 0.5\,\mathrm{m}))$;
\item the first query is the object point nearest to the box centre; the others are drawn with probability
      proportional to the local point density times a weight that decreases with the distance from the centre;
\item the 256 farthest-point queries are kept, and the object queries are added to them.
\end{itemize}
The detector already has a switch that places given 3D points as queries (built for a test with ground-truth
centres), so the object queries can enter through the same path.

\paragraph{Frame and scale.}
The SLAM points are in the SLAM frame and scale; the queries live in the frame of the first input image.
The detector, as it is run today, receives the ground-truth pose of the first frame and a ground-truth scale
(the mean distance of the scan points from the first camera) in training and at test time. Aligning the SLAM
objects to the scene with the true scale, by a similarity fit of the camera centres, is therefore consistent
with the present protocol. A detector that sees the video only would have to use the list's own scale estimate
instead; Table~\ref{tab:arch-cov} shows what that costs.

\paragraph{Check before training: do the queries land in the ground-truth boxes?}
We sampled queries as above for the 88 ScanNet scenes of VSI-Bench and counted the ground-truth boxes of the
detector that contain at least one query (Table~\ref{tab:arch-cov}).

\begin{table}[t]
\centering
\caption{Ground-truth boxes that contain at least one query (88 scenes, 1,273 boxes), by the box's largest side.
Farthest-point: 256 samples of the SLAM cloud. Object queries: from SAM 3 boxes above the given score.
Centre: boxes with a query within 0.1\,m of the box centre.}
\label{tab:arch-cov}
\begin{tabular}{lrrrrrrr}
\toprule
Query set & Queries & All & $<0.3$\,m & 0.3--0.6\,m & 0.6--1.2\,m & $>1.2$\,m & Centre \\
\midrule
\multicolumn{8}{l}{\emph{Aligned with the true scale}} \\
Farthest-point, 256              & 256 & 84.0\% &  3.7\% & 65.8\% & 85.1\% & 93.8\% &  4.8\% \\
Object queries, no threshold     & 362 & 84.3\% & 25.9\% & 80.0\% & 85.5\% & 87.7\% & 16.9\% \\
Object queries, score $\geq 0.3$ & 302 & 82.4\% & 22.2\% & 78.9\% & 83.9\% & 85.3\% & 14.5\% \\
Object queries, score $\geq 0.7$ & 188 & 77.5\% & 14.8\% & 74.2\% & 79.2\% & 80.2\% & 13.6\% \\
Object ($\geq 0.3$) + farthest-point, 256 in total & 256 & 78.2\% & 14.8\% & 73.7\% & 79.2\% & 82.3\% & 13.2\% \\
\midrule
\multicolumn{8}{l}{\emph{Aligned with the list's own scale estimate}} \\
Farthest-point, 256              & 256 & 70.3\% &  0.0\% & 52.1\% & 74.5\% & 76.4\% &  2.4\% \\
Object queries, score $\geq 0.3$ & 302 & 67.6\% & 14.8\% & 61.6\% & 73.0\% & 66.9\% &  9.9\% \\
\bottomrule
\end{tabular}
\end{table}

\begin{itemize}
\item \emph{Small objects gain.} Below 0.6\,m the object queries cover 79--80\% of the boxes against 66\% for
      farthest-point sampling; below 0.3\,m, 22--26\% against 4\% (27 boxes only).
\item \emph{Large objects lose.} Above 1.2\,m farthest-point sampling covers 94\%, object queries 85--88\%,
      because SAM 3 misses some large objects entirely.
\item \emph{Centres are closer.} Three times as many boxes have a query within 0.1\,m of their centre.
\item \emph{Replacing is worse than adding.} With a fixed budget of 256, mixing the two kinds covers fewer boxes
      (78\%) than either kind alone. The object queries must be added to the 256, not take their place.
\item \emph{Scale matters more than the kind of query.} With the estimated scale every row loses about 14
      points.
\end{itemize}

\subsection{Labels from the proposals and a keep decision (task2-det-matching)}
\label{sec:vggtdet-arch-match}

\paragraph{Basic test: the detector as a refiner.}
The label of an object query is the label of its SAM 3 box, and the class head is not used. A new output decides
for each SAM 3 box whether it is kept.
\begin{itemize}
\item Each SAM 3 box is assigned beforehand to at most one ground-truth box of the same label (its centre inside
      the ground-truth box, or the reverse, or IoU $\geq 0.25$; one-to-one, best score first).
\item The queries of an assigned box regress the centre and size of that ground-truth box and have keep $=1$.
      The queries of an unassigned box have keep $=0$ and no box loss. No Hungarian matching is needed for them.
\item At test time the queries of one SAM 3 box are pooled into one output box, so no non-maximum suppression is
      needed.
\end{itemize}
The detector then only refines centre and size and removes extra boxes, which is what the list lacks.

\paragraph{Ceiling of the basic test.}
Table~\ref{tab:arch-keep} gives the result of the assignment on the 88 scenes, for the classes that have
ground truth there.

\begin{table}[t]
\centering
\caption{SAM 3 boxes assigned one-to-one to ground-truth boxes of the same label (88 scenes, true scale).}
\label{tab:arch-keep}
\begin{tabular}{rrrrr}
\toprule
SAM 3 score & Boxes per scene & Keep & Drop & Ground truth with a same-label box \\
\midrule
none       & 31.9 & 31.0\% & 69.0\% & 68.4\% \\
$\geq 0.3$ & 26.0 & 37.2\% & 62.8\% & 66.8\% \\
$\geq 0.7$ & 17.0 & 52.9\% & 47.1\% & 62.4\% \\
\bottomrule
\end{tabular}
\end{table}

\begin{itemize}
\item \emph{Recall is capped at about two thirds.} A detector that only refines and filters cannot find the
      remaining third, and a lower threshold changes little (62\% to 68\%).
\item \emph{The weak classes} are picture (20\% of its boxes contain a same-label query at score $\geq 0.3$),
      trash bin (48\%), cabinet (57\%), window (59\%) and door (67\%); chair, sofa, bed, toilet and desk are at
      89--96\%.
\item \emph{The keep target is not rare.} Between half and two thirds of the proposals are negatives, so the keep
      output has examples of both kinds.
\end{itemize}

\paragraph{Two limits of the basic test.}
It cannot add an object that SAM 3 missed. And a correct SAM 3 box on an object that has no annotation is trained
as ``drop''; this is most likely for ARKitScenes, so the test starts on ScanNet.

\paragraph{Second stage: both kinds of queries.}
The free queries keep the class head and Hungarian matching and can find the objects SAM 3 missed. For the object
queries the label enters in two places: in the matching cost, as a penalty for a ground-truth box of another
label or as a restriction to boxes of the same label; and in the query content, as an embedding of the label
together with the SAM 3 score and the proposal's size, where the content is zero today.

\subsection{Object features from the 2D detections (task2-det-feature)}
\label{sec:vggtdet-arch-feat}

A SAM 3 object is seen as a mask in several frames. Its 2D evidence can be aggregated over those frames into
one feature vector and given to the detector as the content of the object's queries. Three sources are possible
(Table~\ref{tab:arch-feat}).

\begin{table}[t]
\centering
\caption{Sources for a per-object feature, cheapest first.}
\label{tab:arch-feat}
\begin{tabular}{llp{0.42\linewidth}}
\toprule
 & Source & Cost \\
\midrule
A & The detector's own VGGT tokens, averaged under each mask and over the object's views, weighted by the
    SAM 3 score & No new SAM 3 run; the detector must be given the SLAM keyframes as input frames \\
B & SAM 3's own embedding of each detection, averaged over the views & SAM 3 must be run again to store them
    (only masks and scores are stored now) \\
C & SAM 3 image-encoder features pooled under the mask & Run again, and larger files \\
\bottomrule
\end{tabular}
\end{table}

Option A comes first: it needs no new data and tells the decoder which image regions belong to the object.
Option B carries the strongest semantic signal and is added only if A helps.

\subsection{Order of work}

\begin{enumerate}
\item Query coverage and the ceiling of the keep decision, without training
      (done on 88 scenes: Tables~\ref{tab:arch-cov} and~\ref{tab:arch-keep}).
\item The same measurement on the ScanNet training scenes, whose 65-class ground truth includes small objects.
\item Basic test on ScanNet: object queries added to the 256 free queries, label from the proposal, keep output,
      regression targets from the fixed assignment. Validation on the 88 ScanNet scenes of VSI-Bench.
\item Second stage with both kinds of queries.
\item Object features, option A.
\end{enumerate}
Each stage is judged by the list counts (extra and missing boxes) and by the VSI-Bench score of the language
model against the unrefined SAM 3 list.

\paragraph{Limits of the check so far.}
The farthest-point row samples the SLAM cloud as a stand-in for the detector's queries, which are sampled from
VGGT's own cloud for the input frames. The ground truth of the 88 scenes has boxes for 16 classes of large
furniture only, so small objects are under-represented. The SAM 3 lists used the fixed 65-name prompts with
plain ``door'', not ``room door''.
